%   Filename    : chapter_1.tex

\chapter{Introduction}
\label{sec:researchdesc}

\section{Overview of the Current State of Technology}
\label{sec:overview}

Next-generation sequencing (NGS) technologies have advanced genomic research because they allow large volumes of DNA sequences to be analyzed efficiently and at relatively low cost \citep{Metzker2010}. Among these platforms, Illumina sequencing remains widely used for short-read genomic analysis due to its high throughput and base-calling accuracy \citep{Bentley2008,Glenn2011}. The short reads produced by Illumina platforms are commonly used for genome assembly, variant detection, population genetics, and other molecular analyses. In mitochondrial genomics, NGS has made it easier to reconstruct complete mitochondrial genomes, which are valuable for studying species identification, phylogenetic relationships, evolutionary studies, investigations of mitochondrial defects. Because mitochondrial genomes are generally small, circular, and  often repetitive, they are suitable targets for short-read sequencing and assembly workflows \citep{Boore1999,Gray2012,Dierckxsens2017}.

Current mitochondrial assembly workflows commonly rely on bioinformatics tools that reconstruct organelle genomes from raw sequencing reads. Tools such as SPAdes, GetOrganelle, MITObim, and NOVOPlasty have been used to assemble mitochondrial or organellar genomes from short-read datasets \citep{Bankevich2012,Hahn2013,Dierckxsens2017,Jin2020}. These tools support the reconstruction of complete or near-complete mitochondrial sequences, but their performance still depends heavily on the quality of the input reads. Sequencing artifacts introduced before or during library preparation can affect downstream assembly. One persistent challenge is the presence of polymerase chain reaction (PCR)-induced chimeras, which are artificial hybrid sequences formed during amplification \citep{Judo1998,Qiu2001}. In mitochondrial assembly, these chimeric reads may contribute to fragmented contigs, false junctions, or incorrect assembly graph structures, especially because of the structure of mitochondrial genomes \citep{Boore1999,Cameron2014}.

Existing chimera detection approaches are commonly grouped into reference-based and de novo methods. Reference-based approaches compare query sequences against curated databases of known non-chimeric sequences, while de novo approaches infer chimeras from patterns within the dataset itself, often using sequence abundance and compositional similarity \citep{Edgar2011,Edgar2016}. Tools such as UCHIME and UCHIME2 implement these principles and are widely associated with amplicon-based chimera detection \citep{Edgar2011,Edgar2016}. DADA2 also includes a chimera removal step as part of its amplicon sequence variant (ASV) inference pipeline \citep{Callahan2016}. Other approaches, such as CATCh and ChimPipe, extend chimera detection through ensemble machine learning (ML) or identifying chimeric junctions via split-reads and discordant paired-end reads \citep{Mysara2015,Rodriguez2017}. However, many of these tools were designed for microbial amplicon studies, ASV-level analysis, or biological chimera detection, rather than direct read-level detection of PCR-induced chimeras in organellar mitochondrial sequencing data.

At the same time, modern machine learning and deep learning (DL) algorithms have become increasingly relevant in genomic sequence analysis. Feature-based machine learning approaches can use extracted features such as k-mer composition, GC content, alignment quality, soft-clipping patterns, and split-alignment indicators to classify chimeric reads \citep{Vervier2015,Li2018,Sedlazeck2018,Ren2020}. Deep learning methods, on the other hand, can learn patterns directly from encoded nucleotide sequences using representations such as one-hot encoding, k-mer tokenization, and embeddings \citep{He2021}. Neural network a  rchitectures such as one-dimensional convolutional neural network (1D CNN) and Bidirectional Gated Recurrent Unit (BiGRU), have been applied in genomic tasks where local motifs, sequence transitions, and positional dependencies are important \citep{Alipanahi2015,Shen2018,FernandezCastillo2022,Tabane2025}. These developments show that machine learning is a useful approach for detecting complex patterns in genomic data.

This study applies these developments to mitochondrial sequencing data from \textit{Sardinella lemuru}, a small pelagic fish species of economic and ecological importance in Philippine waters \citep{Willette2011,Labrador2021}. Reliable mitochondrial assemblies for this species can support studies on population structure, fisheries genetics, and evolutionary relationships. In the local setting, the study aimed to provide support in  mitochondrial genome assembly workflows done in Philippine Genome Center Visayas (PGC Visayas). To address the need for read-level detection of PCR-induced chimeras, this study presents MitoChime, a machine learning pipeline for identifying chimeric reads in mitochondrial Illumina data. The pipeline analyzed the performance of three different machine learning models, providing a specialized approach for improving the quality and reliability of downstream mitochondrial genome assembly.

\section{Problem Statement}

Chimeric reads can distort assembly graphs and cause misassemblies, with particularly severe effects in mitochondrial genomes \citep{Boore1999,Cameron2014}. Existing assembly pipelines such as SPAdes, GetOrganelle, MITObim, and NOVOPlasty assume that sequencing reads are free of such artifacts \citep{Bankevich2012,Hahn2013,Dierckxsens2017,Jin2020}. At PGC Visayas, several mitochondrial assemblies have failed or yielded incomplete contigs despite sufficient coverage, suggesting that undetected chimeric reads compromise assembly reliability. Meanwhile, previously mentioned chimera detection tools were developed primarily for amplicon-based community analysis and rely heavily on reference or taxonomic comparisons. These approaches are unsuitable for single-species organellar data, where complete reference genomes are often unavailable.

\section{Research Objectives}
\label{sec:researchobjectives}

\subsection{General Objective}
\label{sec:generalobjective}

This study aims to develop and evaluate a machine learning-based pipeline that detects PCR-induced chimeric reads in \textit{Sardinella lemuru} mitochondrial sequencing data in order to improve the quality and reliability of downstream mitochondrial genome assemblies.

\subsection{Specific Objectives}
\label{sec:specificobjectives}

Specifically, the study aims to:
\begin{enumerate}
  \item construct simulated \textit{Sardinella lemuru} Illumina paired-end datasets containing both clean and PCR-induced chimeric reads,
  \item extract read-level alignment-based and sequence-derived features associated with PCR-induced chimerism,
  \item develop, train, and evaluate classical machine learning and neural network models for classifying reads as clean or chimeric,
  \item determine feature importance and identify the informative indicators of PCR-induced chimerism,
  \item identify suitable classifier configurations for pipeline integration by comparing predictive performance, interpretability, runtime, read retention, and downstream assembly behavior, and integrate the selected model options into a modular pipeline deployable on standard computing environments at bioinformatics laboratories, and
  \item assess the effect of the pipeline on downstream mitochondrial genome assembly quality.
\end{enumerate}

\section{Scope and Limitations of the Research}
\label{sec:scopelimitations}

This study focused solely on PCR-induced chimeric reads in \textit{Sardinella lemuru} mitochondrial sequencing data, with the species choice guided by four considerations: (1) to limit interspecific variation in mitochondrial genome size, GC content, and repetitive regions so that differences in read patterns can be attributed more directly to PCR-induced chimerism, (2) to align the analysis with relevant \textit{S. lemuru} sequencing projects at PGC Visayas, (3) to take advantage of the availability of \textit{S. lemuru} mitochondrial assemblies and raw datasets in public repositories such as the National Center for Biotechnology Information (NCBI), which facilitates reference selection and benchmarking, and (4) to develop a tool that directly supports local studies on \textit{S. lemuru} population structure and fisheries management.

The study emphasized \texttt{wgsim}-based simulations and selected empirical mitochondrial datasets from \textit{S. lemuru}. It excludes naturally occurring chimeras, nuclear mitochondrial pseudogenes (NUMTs), and large-scale assembly rearrangements in nuclear genomes. Testing on long-read platforms (e.g., Nanopore, PacBio) and other taxa is outside the scope of this project.

Other limitations in this study include the following: simulations with varying error rates were not performed, so the effect of different sequencing errors on model performance remains unexplored; alternative parameter settings, including k-mer lengths and microhomology window sizes, were not systematically tested, which could affect the sensitivity of both k-mer and microhomology feature detection. Furthermore, machine learning models rely on supervised training with labeled examples, which may limit their ability to detect novel or unexpected chimeric patterns. Finally, hybrid feature-sequence models were not explored, which could potentially capture complementary information for improved classification.

\section{Significance of the Research}
\label{sec:significance}

This research provides both methodological and practical contributions to mitochondrial genomics and bioinformatics. First, MitoChime detects PCR-induced chimeric reads prior to genome assembly, improving the contiguity and correctness of \textit{Sardinella lemuru} mitochondrial assemblies. Second, it replaces informal manual curation with a documented workflow, improving automation and reproducibility. Third, the study contributes to the limited body of work applying machine learning approaches to read-level PCR-induced chimera detection in mitochondrial sequencing data. Fourth, the resulting pipeline is designed to run on computing infrastructures commonly available in regional or resource-constrained laboratories. Finally, more reliable mitochondrial assemblies for \textit{S. lemuru} provide a stronger basis for downstream applications in the field of fisheries and genomics.